%% ASE2026 Project Paper — ACM SIGPLAN format (POPL-style template)
%% Compile on Overleaf: ACM Conference Proceedings Primary Article Template, sigplan
%% Local: see README.md

\documentclass[sigplan,review,anonymous]{acmart}

\setcopyright{none}
\copyrightyear{2026}
\acmYear{2026}
\acmDOI{}
\acmConference[ASE '26]{Advanced Software Engineering Project}{June 2026}{}
\acmISBN{}

\usepackage{booktabs}
\usepackage{graphicx}
\usepackage{listings}
\lstset{basicstyle=\ttfamily\small, breaklines=true}

\begin{document}

\title{Real-Time Interactive Academic Research Assistant with RAG and Citations}

\author{Your Name}
\affiliation{%
  \institution{Your University}
  \city{Your City}
  \country{Your Country}}
\email{your.email@university.edu}

\begin{abstract}
The volume of academic publications has made literature review increasingly time-consuming,
while many AI assistants produce answers without verifiable sources.
This project implements a \emph{Retrieval-Augmented Generation} (RAG) platform that lets
researchers upload PDF and DOCX documents, ask questions in natural language, and receive
\emph{streaming, citation-grounded} responses.
The system supports four pipeline modes---Q\&A, full-document Summarize, Literature search,
and multi-paper Synthesis---and offers \emph{Fast} and \emph{Deep} retrieval settings to
trade latency against thoroughness.
Built with Django REST, Celery, PostgreSQL with pgvector, and a Next.js frontend, the
platform demonstrates how interactive research workflows can combine generative AI with
source traceability.
Evaluation on project \emph{AI in Health Care} (3 indexed PDFs, 193 chunks) shows Fast Q\&A
in approximately 2--4\,s with three citations versus Deep with ten citations at comparable
latency; Literature and Synthesis run in 5--11\,s; full-document Summarize takes 45\,s--4\,min
depending on PDF size.
\end{abstract}

\begin{CCSXML}
<ccs2012>
<concept>
<concept_id>10011036.10011067.10011071</concept_id>
<concept_desc>Computing methodologies~Natural language processing</concept_desc>
<concept_significance>500</concept_significance>
</concept>
<concept>
<concept_id>10002951.10003317</concept_id>
<concept_desc>Information systems~Information retrieval</concept_desc>
<concept_significance>500</concept_significance>
</concept>
<concept>
<concept_id>10003033.10003034</concept_id>
<concept_desc>Human-centered computing~Interactive systems and tools</concept_desc>
<concept_significance>300</concept_significance>
</concept>
</ccs2012>
\end{CCSXML}

\ccsdesc[500]{Computing methodologies~Natural language processing}
\ccsdesc[500]{Information systems~Information retrieval}
\ccsdesc[300]{Human-centered computing~Interactive systems and tools}

\keywords{RAG, academic research assistant, citations, pgvector, streaming LLM, literature review}

\maketitle

\section{Introduction}
\label{sec:intro}

The rapid growth of academic publications has made it increasingly difficult for researchers
and students to efficiently locate, analyze, and synthesize relevant information.
Traditional research methods require extensive manual effort to search for papers, evaluate
their relevance, and extract key insights.
While recent advances in artificial intelligence have introduced tools to assist with these
tasks, many systems lack reliability---often generating responses without verifiable sources
or producing inaccurate information~\cite{ase2026proposal}.

This project addresses these challenges by building an AI-powered academic research assistant
based on a RAG architecture~\cite{lewis2020rag}.
Unlike generic chatbots, the system restricts answers to \emph{user-uploaded documents},
attaches inline citations with page numbers, and optimizes for \emph{real-time interaction}
through server-sent event (SSE) streaming and a Fast mode with reduced retrieval depth.

Most research assistants optimize answer quality at the cost of speed.
Our design explicitly targets \emph{time-to-first-response} and continuous conversation flow:
users receive immediate pipeline feedback, streaming tokens, and a Fast mode for low-latency
interaction while reading.

\paragraph{Contributions.}
To summarize, our contributions are:
\begin{itemize}
  \item A multi-mode RAG platform with four specialized pipelines: Q\&A, Summarize,
        Literature search, and Synthesis.
  \item A Fast vs.\ Deep retrieval strategy with measurable latency--quality tradeoffs.
  \item A streaming chat UX with pipeline stage feedback and time-to-first-token metrics.
  \item Citation precision with per-chunk snippets, character offsets, and click-to-highlight
        excerpt panels.
  \item A two-phase \emph{Refine with Deep} pattern for upgrading Fast Q\&A drafts on demand.
\end{itemize}

\paragraph{Paper structure.}
\S\ref{sec:problem} states the problem and requirements.
\S\ref{sec:arch} describes the system architecture.
\S\ref{sec:rag} presents the RAG pipeline and ingestion methodology.
\S\ref{sec:modes} surveys the four pipeline modes.
\S\ref{sec:impl} discusses implementation highlights.
\S\ref{sec:eval} reports experimental evaluation.
\S\ref{sec:related} surveys related work.
\S\ref{sec:conclusion} concludes.

\section{Problem Statement and Requirements}
\label{sec:problem}

Researchers face three interrelated problems: (1)~\emph{information overload}---too many
papers to read in available time; (2)~\emph{untrusted AI outputs}---hallucinated facts and
missing citations limit academic use; and (3)~\emph{slow AI UX}---systems that prioritize
answer perfection over interaction speed disrupt research flow.

\paragraph{Functional requirements.}
The system must support user registration and JWT authentication; project-scoped document
upload (PDF, DOCX); async ingestion (extract, chunk, embed, index); chat with streaming
responses and history; automatic inline citations; and multiple research modes.

\paragraph{Non-functional requirements.}
Fast Q\&A on warm queries should complete in under 10\,s; answers must be traceable to
source document, page, and excerpt; the backend must scale via worker queues; and embeddings
must run on CPU without requiring a GPU.

Table~\ref{tab:proposal} maps the original project proposal objectives to our implementation.

\begin{table}[t]
  \caption{Proposal objectives vs.\ implementation.}
  \label{tab:proposal}
  \small
  \begin{tabular}{@{}p{0.42\columnwidth}p{0.48\columnwidth}@{}}
    \toprule
    \textbf{Objective} & \textbf{Implementation} \\
    \midrule
    Reduce research time & Fast mode ($k{=}3$), SSE streaming, Celery ingestion \\
    Improve accuracy & Project-scoped RAG; no open-web hallucination \\
    Ensure scalability & Docker Compose; Celery workers; HNSW pgvector index \\
    Enhance trust & Inline citations; Sources panel; excerpt highlight API \\
    \bottomrule
  \end{tabular}
\end{table}

\section{System Architecture}
\label{sec:arch}

Figure~\ref{fig:arch} summarizes the major components.
The \textbf{Next.js frontend} provides authentication, project workspace, document list,
streaming chat, and a citation panel.
The \textbf{Django REST API} exposes JWT auth, project/document CRUD, and chat/SSE endpoints.
\textbf{Celery workers} handle background ingestion and full-document summarize jobs.
\textbf{PostgreSQL with pgvector} stores relational data and 384-dimensional embeddings
with an HNSW cosine index~\cite{pgvector2023,johnson2019hnsw}.
\textbf{Redis} serves as the Celery broker.
\textbf{sentence-transformers} computes local CPU embeddings (\texttt{all-MiniLM-L6-v2}).
The \textbf{OpenAI Chat API} generates answers (\texttt{gpt-4o-mini}).

\begin{figure}[t]
  \centering
  \fbox{\parbox{0.92\columnwidth}{\small
    \textbf{Next.js UI} $\leftrightarrow$ \textbf{Django REST} (JWT + SSE)\\
    $\downarrow$ \hspace{2em} $\downarrow$ \hspace{2em} $\downarrow$\\
    PostgreSQL+pgvector \quad Redis/Celery \quad OpenAI Chat API\\
    \textit{Ingestion: PDF/DOCX $\to$ chunk $\to$ embed $\to$ HNSW index}
  }}
  \caption{High-level system architecture.}
  \label{fig:arch}
\end{figure}

\paragraph{Data model.}
A \emph{Project} contains documents and chat messages with per-user isolation.
A \emph{Document} tracks upload status (\texttt{pending} $\to$ \texttt{processing} $\to$
\texttt{ready} / \texttt{failed}).
A \emph{DocumentChunk} stores text, embedding, page number, and metadata.
\emph{ChatMessage} records user/assistant turns with citations JSON and mode metadata.

\paragraph{Security.}
All project routes require JWT authentication and project membership.
Retrieval queries filter by \texttt{project\_id} and \texttt{document.status = ready}.

\section{RAG Pipeline}
\label{sec:rag}

\subsection{Ingestion}
Users upload PDF or DOCX files (max 50\,MB).
A Celery task extracts text per page, splits it into chunks of 300--500 tokens, embeds each
chunk with \texttt{all-MiniLM-L6-v2}~\cite{reimers2019sbert} (normalized 384-dim vectors),
and stores rows in \texttt{DocumentChunk} with an HNSW index for cosine similarity search.

\subsection{Query Processing (Q\&A Mode)}
The Q\&A pipeline executes seven stages:
(1)~load recent chat history (Fast: 6 turns, Deep: 20);
(2)~enhance and embed the user query;
(3)~retrieve top-$k$ chunks via pgvector cosine search, project-scoped;
(4)~aggregate excerpts grouped by document with page labels;
(5)~generate an answer via OpenAI chat with system prompt and history;
(6)~build citations from retrieved chunks (snippet, offsets, \texttt{chunk\_id});
(7)~persist user and assistant messages.
Overview-style queries merge intro and early-page chunks with vector results to improve recall.

\subsection{Prompt Design}
The Q\&A system prompt instructs the model to answer only from retrieved excerpts and
conversation history; provide partial answers when excerpts are relevant but incomplete;
respond \emph{Not found in the provided documents} only when context truly lacks support;
and add inline citations in the form \texttt{[DocumentName -- Page N]}.

\section{Pipeline Modes}
\label{sec:modes}

Table~\ref{tab:modes} summarizes the four pipeline modes and recommended use cases.

\begin{table}[t]
  \caption{Pipeline modes and retrieval settings.}
  \label{tab:modes}
  \small
  \begin{tabular}{@{}llcc@{}}
    \toprule
    \textbf{Mode} & \textbf{Best for} & \textbf{Fast $k$} & \textbf{Deep $k$} \\
    \midrule
    Q\&A & Specific factual questions & 3 & 10 \\
    Summarize & Full overview of one doc & \multicolumn{2}{c}{all chunks} \\
    Literature & Cross-paper exploration & 15 & 25 \\
    Synthesis & Multi-paper lit.\ review & 20 & 40 \\
    \bottomrule
  \end{tabular}
\end{table}

\paragraph{Q\&A.}
Answers specific questions using top-$k$ vector retrieval with inline citations.

\paragraph{Summarize.}
Hierarchical map-reduce over \emph{all} chunks of one selected document, producing sections:
Objective, Methodology, Key Findings, Limitations, Conclusion.
Runs asynchronously via Celery.

\paragraph{Literature search.}
Exploration report with Thematic clusters, Paper-level summaries, and Key insights.

\paragraph{Synthesis.}
Literature-review style output with themed sections, Agreements, Contradictions, and Conclusion.

\section{Implementation Highlights}
\label{sec:impl}

\paragraph{Streaming chat (SSE).}
\texttt{POST /api/projects/\{id\}/chat/stream/} returns \texttt{text/event-stream} events:
\texttt{status} (pipeline stage), \texttt{token} (answer delta), and \texttt{done}
(full answer, citations, \texttt{message\_id}, \texttt{latency\_ms}).
The frontend uses \texttt{fetch} + \texttt{ReadableStream} to support JWT headers.

\paragraph{Fast vs.\ Deep.}
Reranking was intentionally removed to reduce latency and API cost.
Users can trigger \emph{Refine with Deep} after a Fast Q\&A answer to regenerate with
deeper retrieval without re-typing the question.

\paragraph{Citation precision.}
Each citation includes document name, page, chunk ID, snippet text, and character offsets.
The frontend fetches full chunk text via
\texttt{GET /api/projects/\{id\}/chunks/\{chunk\_id\}/} for highlight rendering.

\section{Evaluation}
\label{sec:eval}

We benchmarked the system on project \emph{AI in Health Care} (3 indexed PDFs, 193 chunks)
using queries from saved chat history and warm \texttt{/chat/benchmark/} runs.

Table~\ref{tab:qa} reports Q\&A results.
Table~\ref{tab:lit} reports Literature search.
Table~\ref{tab:syn} reports Synthesis.

\begin{table}[t]
  \caption{Q\&A (healthcare chat queries).}
  \label{tab:qa}
  \small
  \begin{tabular}{@{}lrrrr@{}}
    \toprule
    \textbf{Query} & \textbf{Fast ms} & \textbf{Deep ms} & \textbf{Fast cite.} & \textbf{Deep cite.} \\
    \midrule
    Problem statement? & 2{,}005 & 2{,}249 & 3 & 10 \\
    What are the Solutions? & 2{,}005 & 4{,}400 & 3 & 10 \\
    Methodology described? & 3{,}376 & 3{,}803 & 3 & 10 \\
    \bottomrule
  \end{tabular}
\end{table}

\begin{table}[t]
  \caption{Literature search latency.}
  \label{tab:lit}
  \small
  \begin{tabular}{@{}lrrrr@{}}
    \toprule
    \textbf{Query} & \textbf{Fast ms} & \textbf{Deep ms} & \textbf{Fast cite.} & \textbf{Deep cite.} \\
    \midrule
    Themes across papers? & 5{,}047 & 6{,}644 & 15 & 25 \\
    Main findings? & 8{,}564 & 10{,}654 & 15 & 25 \\
    \bottomrule
  \end{tabular}
\end{table}

\begin{table}[t]
  \caption{Synthesis latency.}
  \label{tab:syn}
  \small
  \begin{tabular}{@{}lrrrr@{}}
    \toprule
    \textbf{Query} & \textbf{Fast ms} & \textbf{Deep ms} & \textbf{Fast cite.} & \textbf{Deep cite.} \\
    \midrule
    Methodology agreements? & 6{,}225 & 7{,}342 & 20 & 40 \\
    \bottomrule
  \end{tabular}
\end{table}

\paragraph{Findings.}
Fast Q\&A delivers interactive responses (2--4\,s warm) with three citations.
Deep Q\&A retrieves $3\times$ more chunks (10 vs 3 citations) at similar latency.
Literature and Synthesis Deep modes add citations (25 vs 15, 40 vs 20) with modest latency
increase on three papers.
Summarize on the 17-chunk EPF brief took 45\,s; larger PDFs scale to several minutes.
First query after backend restart incurs a cold start for embedding model load;
startup warm-up mitigates this effect.

Qualitative trust checks confirm that citations map to real uploaded pages, unsupported
questions return the exact not-found phrase, and citation clicks open highlighted excerpts.

\section{Related Work}
\label{sec:related}

\paragraph{RAG.}
Lewis et al.~\cite{lewis2020rag} combine dense retrieval with LLMs to ground generation in
external documents.
We apply RAG in a project-scoped academic setting rather than open-web search.

\paragraph{Academic research tools.}
Semantic Scholar, Elicit, and Connected Papers focus on discovery across published corpora.
Our system complements these by operating on a private document library per project.

\paragraph{Citation-grounded QA.}
Systems such as Perplexity provide source links but search the public web.
We enforce grounding through prompt design, retrieval scope, and explicit not-found responses.

\paragraph{Interactive LLM applications.}
Streaming improves perceived latency; our design targets time-to-first-response as a
first-class metric for research chat workflows.

\section{Conclusion}
\label{sec:conclusion}

This project delivers an AI-powered academic research assistant that fulfills and extends
the original ASE2026 proposal.
By combining RAG with streaming interaction, four specialized pipeline modes, and verifiable
citations, the system demonstrates how generative AI can support academic workflows without
sacrificing trust or responsiveness.
The Fast/Deep tradeoff provides a practical pattern for balancing latency and thoroughness.

Future work includes production hardening, a full PDF page viewer, formal citation-precision
evaluation, hybrid BM25+vector retrieval, and optional reranking in Deep mode only.

\bibliographystyle{ACM-Reference-Format}
\bibliography{references}

\end{document}
